\chapter{Bleeding Control}

\section*{Definition and Overview}
Bleeding control is the set of first-aid measures used to stop or slow severe bleeding while emergency help is on its way. External bleeding occurs when a blood vessel is cut or torn and blood escapes through a break in the skin. Most bleeding from small cuts and grazes stops quickly with simple measures, but bleeding from a major blood vessel can be life-threatening within minutes. The aims of bleeding control are to slow blood loss, protect the wound from contamination, prevent shock, and arrange transfer to emergency care as quickly as possible.

\section*{Epidemiology}
Severe bleeding is one of the leading preventable causes of death after injury worldwide, whether from road traffic crashes, falls, workplace accidents, or penetrating injuries. Uncontrolled haemorrhage can cause shock within minutes, and the outcome is strongly time-dependent: early, effective first aid measurably improves survival. Anyone may be involved in an accident, so basic bleeding-control skills are valuable for the whole community, and they are especially important in rural or remote settings where professional help may be delayed.

\section*{Aetiology and Pathophysiology}
Severe bleeding most often follows trauma, but the ability to control it also depends on the person's own clotting mechanisms:

\begin{itemize}
    \item \textbf{Cuts and lacerations} from sharp objects, tools, machinery, or glass
    \item \textbf{Crush injuries} that rupture blood vessels beneath intact or damaged skin
    \item \textbf{Fractures} that tear adjacent vessels, for example a broken leg or upper arm
    \item \textbf{Penetrating injuries} from knives, bullets, or other objects
    \item \textbf{Anticoagulant medication or bleeding disorders}, which impair the body's natural clotting and make bleeding harder to control
\end{itemize}

When a vessel is damaged, the body responds by constricting the vessel and forming a clot. Severe or rapid blood loss overwhelms these defences: the circulating blood volume falls, blood pressure drops, and vital organs receive less oxygen, producing the state known as haemorrhagic (hypovolaemic) shock. Without intervention, shock progresses to organ failure and death.

\section*{Clinical Features}
The outward features of severe bleeding differ from those of minor bleeding, and recognising them triggers urgent action:

\begin{itemize}
    \item Blood pooling on the ground or soaking through clothing rapidly
    \item Bleeding that does not slow despite firm pressure
    \item Pale, cool, or clammy skin
    \item A rapid, weak pulse
    \item Rapid, shallow breathing
    \item Dizziness, weakness, confusion, or agitation
    \item Fainting, or a feeling of impending collapse
    \item Internal bleeding may present with shock and no visible blood, for example after major trauma
\end{itemize}

\section*{Differential Diagnosis}
Identifying the type of bleeding gives clues to the vessels involved and guides management:

\begin{itemize}
    \item \textbf{Arterial bleeding:} bright red blood that spurts with each heartbeat and is the most rapidly life-threatening
    \item \textbf{Venous bleeding:} darker red blood that flows steadily but usually with less force
    \item \textbf{Capillary bleeding:} a slow ooze from abrasions and grazes that generally stops with light pressure
    \item \textbf{Internal bleeding:} no external blood but features of shock after an injury
    \item \textbf{Bleeding from natural openings:} the nose, ears, mouth, bowel, or vagina, suggesting injury or underlying disease
    \item \textbf{Minor but alarming bleeding:} such as nosebleeds, which look dramatic but are usually controllable with simple measures
\end{itemize}

\section*{When to Seek Medical Care}
All wounds other than the most trivial need medical review. Seek professional assessment for wounds that are deep, gaping, or dirty, that involve the face, hands, joints, or genitals, that contain a foreign body, or where bleeding has not stopped.

\textbf{Call 000 (or local emergency services) for:}
\begin{itemize}
    \item Bleeding that soaks through dressings despite firm, sustained pressure
    \item Bleeding from a major wound, especially a partial or complete amputation
    \item Signs of shock: a pale, clammy, confused, drowsy, or barely conscious person
    \item Any bleeding that cannot be controlled within a few minutes
    \item Bleeding associated with a severe injury or a suspected broken bone
\end{itemize}

\section*{Diagnosis and Investigations}
In the emergency setting, control of bleeding takes priority, and investigation follows once the patient is stable:

\begin{itemize}
    \item \textbf{Wound assessment:} careful examination of the size, depth, and involvement of vessels, nerves, tendons, or bone
    \item \textbf{Blood tests:} full blood count, clotting studies, and cross-matching if transfusion is required
    \item \textbf{Imaging:} X-ray and CT scanning to identify internal injuries, fractures, or retained foreign bodies
    \item \textbf{Monitoring:} continuous measurement of blood pressure, heart rate, oxygen saturation, and urine output to detect shock
    \item \textbf{Scans for internal bleeding:} focused ultrasound or CT in abdominal and chest trauma
\end{itemize}

\section*{Management}
The sequence of first aid for severe external bleeding is simple and life-saving:

\begin{itemize}
    \item \textbf{Direct pressure:} firm, sustained pressure on the wound with a clean cloth or pad
    \item \textbf{Bandaging:} once bleeding slows, secure a firm bandage over the pressure dressing
    \item \textbf{Do not remove} a blood-soaked dressing -- add another layer on top and continue pressing
    \item \textbf{Pressure points:} if direct pressure fails, press firmly on the artery above the wound (for example in the groin or armpit)
    \item \textbf{Tourniquet:} for life-threatening limb bleeding that pressure cannot control; apply high and tight, note the time, and do not remove it until medical help arrives
    \item \textbf{Position:} lay the person down, keep them warm and still, and raise their legs if this is comfortable
    \item \textbf{Emergency transfer:} call an ambulance immediately, or transport the person to hospital if it is quicker
    \item \textbf{Hospital treatment:} surgical exploration, vessel repair, stitches, or blood transfusion as required
\end{itemize}

\section*{Complications}
\begin{itemize}
    \item Hypovolaemic shock from rapid blood loss
    \item Wound infection from contamination
    \item Damage to nerves, tendons, bones, or major vessels from the injury itself
    \item Ischaemic tissue damage if a tourniquet is left on for too long
    \item Anaemia after significant blood loss
    \item Scarring, and psychological distress following a serious injury
    \item Late effects of massive transfusion, including clotting abnormalities
\end{itemize}

\section*{Prognosis and Outcomes}
The outcome depends on the site and severity of the bleeding and the speed with which definitive care is reached. With prompt bleeding control, fluid resuscitation, and surgical treatment, most people who reach hospital with severe external bleeding survive. Rapid and effective first aid before arrival is among the strongest predictors of a good outcome after major injury, which is why public training in bleeding control saves lives.

\section*{Prevention and Self-Care}
\begin{itemize}
    \item Learn basic first aid and bleeding-control skills through an accredited course, and refresh them regularly
    \item Keep a first-aid kit with sterile dressings, bandages, and gloves at home, at work, and in the car
    \item Use protective equipment and safe techniques with knives, tools, and machinery
    \item Supervise children around sharp objects and tools
    \item Treat minor cuts promptly: clean the wound, apply pressure, cover it, and keep it dry
    \item If you take blood-thinning medication, mention it to any health professional and seek advice about seemingly minor wounds
\end{itemize}

\section*{Sources and Further Reading}
The following sources provide reliable, up-to-date information for patients, first responders, and students of medicine.

\begin{itemize}
    \item NHS. \textit{How to stop bleeding -- first aid guidance.} https://www.nhs.uk/conditions/first-aid/how-to-stop-bleeding/
    \item World Health Organization. \textit{Global guidance on trauma and emergency care.} https://www.who.int/
    \item MedlinePlus. \textit{Bleeding and first aid for bleeding -- patient information.} https://medlineplus.gov/
    \item MSD Manual. \textit{Emergency procedures for external bleeding -- professional edition.} https://www.msdmanuals.com/
    \item Mayo Clinic. \textit{First aid for bleeding wounds.} https://www.mayoclinic.org/first-aid/
\end{itemize}
